\begin{minipage}{0.6\textwidth}
\subsection*{Drill: Defence in Depth}
\label{drill:three_wall/defence_in_depth}

This drill requires a three-wall and a jammer.
This drill is about using the depth of the wall rather than trying to keep the jammer on the butts. 

Start with the jammer two feet behind the wall.  
The wall should be set up such that the butts are not touching the brace - similar to the no-hands drill.


\end{minipage}
\hfill
\hspace{0.5cm}
\begin{minipage}{0.38\textwidth}
\begin{center}
\includegraphics[scale=0.8,trim={0 0 0 4cm}, clip]{img/three_wall/depth/line}

{\it The brace has a shorter path distance to cover than the jammer.} 
\end{center}
\end{minipage}
\vspace{0.5cm}


The goal is to use the butts as a screen to force a jammer commitment. 
The jammer will have to take an additional second or so to move around the butts, giving the brace time to cover the jammer's line and slightly impede the jammer. 
Using the time from this impediment the butts should then close on the jammer, pocketing. 
To achieve this, the brace needs to be free to move and actively track the jammer at all times - without being impeded by contact with their wall.
Simultaneously, to prevent a jammer cutting back through the middle, the brace needs to be close enough to the wall to cover that seam.

The butt on the side that the jammer is taking needs to commit to taking the jammer out - if they give up early then the jammer will have sufficient space to juke the brace.   

The jammer should focus on having sufficient space to juke the brace, and if the trap closes then to be able to pull back and take the other line.


